\scanpage{35}
\uline{Définition} 2.5. Soient $X,Y,\ldots$ et considérons une
application
\[
\varphi:\underline{Y(k)}\longrightarrow \tikz[baseline=(terme.base)]{\node[inner sep=1pt] (terme) {$A^i_{\mathrm{alb}}(X)$};\draw[gray,thin,rotate around={-5:(terme.center)}] (terme.center) ellipse (2.7em and .85em);}
\]
On dit que cette application est \uline{algébrique} si

$1^\circ)$ L'hom $\deleted{Z}\operatorname{Cyc}_0(Y)\longrightarrow A^i_{\mathrm{alb}}(X)$
correspondant se factorise par
\margin{\diagram{\begin{tikzcd}[ampersand replacement=\&,column sep=tiny,font=\scriptsize]
Y\arrow[r]\& A_Y^1\arrow[r,"N"]\& A_Y^N
\end{tikzcd}}}
$\deleted{\unsure{A_Y}}A_{0\,\mathrm{alb}}(Y)\simeq A_Y^*(k)$, i.e. \deleted{\unsure{que|s'annule}} est
nulle sur les $\unsure{\mathfrak{a}|\mathfrak{z}|\mathfrak{t}}$ de degré 0 tels
que $S_Y(\mathfrak{a})=0$

$2^\circ)$ L'hom induit
\[
\varphi:\tikz[baseline=(terme.base)]{\node[inner sep=1pt] (terme) {$A_Y^{\unsure{0|*}}(k)$};\draw[gray,thin,rotate around={-5:(terme.center)}] (terme.center) ellipse (2.7em and .85em);}\longrightarrow A^i_{\mathrm{alb}}(X)
\]
est tel que pour $N>0$ convenable, $\underline{N\varphi}(x)=\varphi(Nx)$
puisse être induit par un cycle algébrique
\[
Z\in A^i_{\mathrm{alb}}(X\times A_Y)
\]

\uline{Prop.} 2.6. Si $\varphi$ est définie par un $Z\in A^i_{\mathrm{alb}}(X\times Y)$,
elle est \uline{algébrique}.

C'est 2.1. Va être généralisé en 2.7. (ii)

\uline{Prop.} 2.7. \inserted{(0) Forment un \uline{ss-groupe}. Si $N\varphi$ est alg., $\varphi$ l'est.}
(i) Fonctorialité en $Y$. \deleted{Plus précisément}

(ii) Fonctorialité en $X$, plus précis, si on a
un hom $\deleted{P}A^i_{\mathrm{alb}}(X)\longrightarrow\deleted{P}A^{\unsure{j|i'}}_{\mathrm{alb}}(X')$ défini par
une classe de correspondance algébrique \uline{à coeff.
\deleted{\unsure{rationnels}}} entiers.
% archive sheet number: 34
